% ===================================================================
%  TAVOLA N. 2 — Il corpo umano. --- Ricreazione. --- Giochi.
%  Traduction 2026 d'apres texto/io, controlee sur texto/fr et
%  texto/en. Consigne : voir texto/it/10-tavola-01.tex.
%
%  LA PREMIERE MOITIE DE CETTE PAGE NOMME LE CORPS ENTIER AVEC DES
%  DIMINUTIFS, ET PERSONNE NE S'EN APERCOIT PLUS.
%
%  Le latin classique disait « oculus », « auris », « genu »,
%  « cerebrum ». Le latin parle a prefere les diminutifs — oculus a
%  cede a « ocellus », auris a « auricula », genu a « genuculum »,
%  cerebrum a « cerebellum » — et c'est cette forme-la, et elle seule,
%  que l'italien a gardee :
%     occhio    < oculus, mais par la forme diminuee
%     orecchio  < auricula   = la petite oreille
%     ginocchio < genuculum  = le petit genou
%     cervello  < cerebellum = le petit cerveau
%     polpaccio < pulpaticium
%  L'ALINEA 4 DECRIT DONC UN VISAGE FAIT DE PETITES CHOSES, et le
%  lecteur italien n'entend plus rien de ce diminutif : « occhio » est
%  pour lui l'oeil, non le petit oeil.
%
%  ET « CERVELLO » A UNE SUITE. Le « cerebellum » latin, qui n'etait
%  qu'un cerveau en petit, est devenu en anatomie moderne un ORGANE A
%  PART — le cervelet, sous le cerveau proprement dit. L'italien
%  courant appelle donc le cerveau du nom que la medecine reserve a
%  autre chose. La colonne ecrit « cervello » a l'alinea 2, parce que
%  c'est le mot de tout le monde ; la medecine dirait « encefalo ».
%
%  LA SECONDE MOITIE, LES JEUX, REPOND AU TABLEAU 2 ALLEMAND.
%
%  Jahn avait refuse « Gymnastik » en 1811 parce que le mot etait grec
%  et passait par le francais, et fabrique « Turnen ». L'italien dit
%  ginnastica, sans y avoir jamais pense. Il n'a pas eu de
%  Jahn, et il n'en a pas eu besoin : ON NE PROUVE PAS QUE LE GYMNASE
%  EST DE CHEZ SOI QUAND IL EST DE CHEZ SOI.
%
%  MAIS L'ALINEA 30 DONNE UN MECANISME QUE LA COLONNE ALLEMANDE
%  N'AVAIT PAS TROUVE, ET C'EST LE SIXIEME.
%
%  Le football arrive d'Angleterre vers 1890. L'allemand le calque —
%  Fußball ; l'estonien aussi — jalgpall ; le romanche aussi —
%  ballape. L'italien, lui, ne calque rien et n'emprunte rien : il
%  l'appelle calcio, qui est le nom d'un jeu florentin du
%  XVIe siecle ou l'on frappait deja une balle du pied. LE MOT N'EST
%  NI EMPRUNTE, NI FORGE, NI CALQUE : IL EST DEPLACE. Une chose neuve
%  recoit le nom d'une chose ancienne qui lui ressemble, et le nom
%  ancien change simplement de titulaire.
%
%  La colonne allemande avait releve cinq facons pour un mot de
%  franchir une frontiere ; en voici une sixieme, et c'est la seule
%  qui ne franchisse aucune frontiere du tout. Elle demande une
%  condition que les cinq autres n'exigent pas : QU'ON AIT DEJA EU LA
%  CHOSE, OU QUELQUE CHOSE D'ASSEZ PROCHE. Le tableau 3 avait dit
%  qu'on ne peut pas reprendre le nom d'une chose qui est la ; ici
%  l'on voit l'inverse — on peut donner le nom d'une chose qui n'y est
%  plus.
% ===================================================================

%%K t02-apar-1 apar
\VUsaut{4.0mm}
\VUtitre{15.0pt}{60}{TAVOLA N. 2}
\VUsaut{2.0mm}
\VUtitre{11.4pt}{0}{Il corpo umano. --- Ricreazione.}
\VUtitre{11.4pt}{0}{Giochi.}

%%K t02-tit-0 sub
\VUtitre{13.2pt}{0}{Il corpo umano.}
\VUsaut{2.4mm}
\VUcentre{10.2pt}{0}{\textit{(Una semplice lezione di scienze naturali.)}}
\VUsaut{3.0mm}

%%K t02-01-1 p
1. --- Le parti principali del corpo umano sono: la \VUgras{testa}, il
\VUgras{tronco} e gli \VUgras{arti}.

%%K t02-tit-1 sub
\VUpk{10.2pt}{40}{La testa.}
\VUsaut{3.0mm}

%%K t02-02-1 p
2. --- La testa ha due parti: il \VUgras{cranio}, che contiene il \VUgras{cervello}
ed è coperto dai \VUgras{capelli}, e il \VUgras{viso}.

%%K t02-03-1 p
3. --- Nel nostro disegno non vediamo né il cranio né il cervello; vediamo però molto
chiaramente le parti importanti del viso.

%%K t02-04-1 p
4. --- Per prima cosa c'è la \VUgras{fronte}, che fa pensare a un'intelligenza viva e
a una mente sempre al lavoro. Le \VUgras{tempie} sono molto scoperte. L'\VUgras{occhio}
è vivace e ben aperto. Un folto \VUgras{sopracciglio} e lunghe \VUgras{ciglia} lo
proteggono.

%%K t02-05-1 p
5. --- Poi vengono il \VUgras{naso} e le \VUgras{narici}. Con il naso odoriamo e
respiriamo. Ai due lati del naso ci sono le \VUgras{guance}, e più in fuori le
\VUgras{orecchie}. La parte inferiore del viso è formata dalla \VUgras{bocca} e dal
\VUgras{mento}.

%%K t02-06-1 p
6. --- Non si vedono molto bene, perché i \VUgras{baffi} e la \VUgras{barba} li
nascondono. Sappiamo però che la bocca ha le due \VUgras{labbra}, la
\VUgras{mascella superiore} e la \VUgras{mascella inferiore} con i loro
\VUgras{denti}, e la \VUgras{lingua}. Con la bocca parliamo e mangiamo. Con la lingua
sentiamo il sapore dei cibi.

%%K t02-07-1 p
7. --- L'occhio, l'orecchio, il naso e la bocca sono, insieme alle dita, gli organi
dei cinque sensi: la \VUgras{vista}, l'\VUgras{udito}, l'\VUgras{olfatto}, il
\VUgras{gusto} e il \VUgras{tatto}.

%%K t02-08-1 p
8. --- La testa è dunque una delle parti più interessanti del corpo umano.

%%K t02-tit-2 sub
\VUsaut{4.4mm}
\VUpk{10.2pt}{40}{Il tronco.}
\VUsaut{3.0mm}

%%K t02-09-1 p
9. --- Il \VUgras{collo} unisce la testa al tronco. Comprende la \VUgras{gola} e la
\VUgras{nuca}.

%%K t02-10-1 p
10. --- Anche il tronco contiene molti organi essenziali. Nel \VUgras{petto} ci sono i
\VUgras{polmoni}, con cui respiriamo, e il \VUgras{cuore}, che è il centro della vita.
Quando il cuore smette di battere, l'essere vivente muore.

%%K t02-11-1 p
11. --- Le \VUgras{costole} sostengono il petto.

%%K t02-12-1 p
12. --- Le altre parti del tronco sono il \VUgras{fianco}, la \VUgras{schiena}, le
\VUgras{spalle} e l'\VUgras{addome} (la pancia). Per dormire ci mettiamo su un fianco
o sulla schiena, e i pesi li portiamo sulla spalla.

%%K t02-13-1 p
13. --- Nell'addome ci sono lo \VUgras{stomaco} e l'\VUgras{intestino}. Sono loro che
digeriscono i nostri cibi.

%%K t02-tit-3 sub
\VUsaut{4.4mm}
\VUpk{10.2pt}{40}{Gli arti.}
\VUsaut{3.0mm}

%%K t02-14-1 p
14. --- Abbiamo quattro \VUgras{arti}: due \VUgras{braccia} e due \VUgras{gambe}. Con
le braccia prendiamo, teniamo, solleviamo e raccogliamo gli oggetti; con le gambe
stiamo in piedi, camminiamo e corriamo.

%%K t02-15-1 p
15. --- La \VUgras{spalla} unisce il braccio al corpo. Un \VUgras{muscolo} molto
robusto, il bicipite, muove il braccio, che il \VUgras{gomito} unisce
all'\VUgras{avambraccio}. Il \VUgras{polso} forma l'articolazione della
\VUgras{mano}, che finisce con le \VUgras{dita} e le \VUgras{unghie}. Sulla mano si
distingue una \VUgras{vena} (e un'arteria) in cui scorre il \VUgras{sangue}.

%%K t02-16-1 p
16. --- Le parti della gamba sono: la \VUgras{coscia}, il \VUgras{ginocchio} e
l'\VUgras{incavo del ginocchio}, il \VUgras{polpaccio}, la cui \VUgras{carne} è
coperta dalla \VUgras{pelle}, la \VUgras{caviglia} e il \VUgras{piede}. Le
\VUgras{ossa} della gamba sono molto grosse e molto forti. All'estremità del piede ci
sono le \VUgras{dita}. Le altre parti sono la \VUgras{pianta} e il \VUgras{tallone}.

%%K t02-17-1 p
17. --- Questa è una descrizione quasi completa del corpo umano.

%%K t02-17-2 p
\textit{Nota.} --- \textit{Con la frase «mi fa male il... (la testa, e così via)» il
maestro può ripassare tutto questo vocabolario dal punto di vista della buona o
cattiva} \VUgras{salute}.

%%K t02-tit-4 sub
\VUsaut{5.0mm}
\VUpk{10.2pt}{40}{Ginnastica.}
\VUsaut{2.4mm}
\VUcentre{10.2pt}{0}{\textit{(Raccontata da uno degli alunni.)}}
\VUsaut{3.0mm}

%%K t02-18-1 p
18. --- Per restare agili, svelti e sani dobbiamo esercitare le gambe e le braccia.
Per questo giochiamo e facciamo \VUgras{ginnastica}.

%%K t02-19-1 p
19. --- Durante la settimana questi due esercizi si fanno nel cortile della scuola
(vedi la tavola n. 1). Ma il giovedì e la domenica andiamo su un grande prato, dove
abbiamo spazio per correre e divertirci.

%%K t02-20-1 p
20. --- I nostri maestri e i nostri genitori vengono con noi qualche volta; ma è il
\VUgras{maestro di ginnastica} \textsuperscript{(12)} che dirige gli esercizi.

%%K t02-21-1 p
21. --- Sul prato, a sinistra dell'ingresso, ci sono gli \VUgras{attrezzi ginnici}
\textsuperscript{(1)}: saliamo sulla \VUgras{scala} \textsuperscript{(2)}, ci
dondoliamo a testa in giù sugli \VUgras{anelli} \textsuperscript{(3)} e sul
\VUgras{trapezio} \textsuperscript{(4)}, ci tiriamo su a forza di braccia fino in cima
alla \VUgras{scala di corda} \textsuperscript{(5)} o alla \VUgras{fune con nodi}
\textsuperscript{(6)}.

%%K t02-22-1 p
22. --- Intanto i nostri amici saltano il \VUgras{cavallo} \textsuperscript{(7)} o le
\VUgras{parallele} \textsuperscript{(11)}. Altri fanno \VUgras{pugilato}
\textsuperscript{(8)} o tirano di \VUgras{scherma} \textsuperscript{(9)} con la
maschera e il \VUgras{fioretto}. Altri ancora sollevano \VUgras{manubri}
\textsuperscript{(10)}.

%%K t02-tit-5 sub
\VUsaut{4.6mm}
\VUpk{10.2pt}{40}{Giochi.}
\VUsaut{3.0mm}

%%K t02-23-1 p
23. --- Intorno ai ginnasti, ragazzi e ragazze fanno i \VUgras{giochi} più diversi. Le
ragazze giocano con il loro \VUgras{cerchio} \textsuperscript{(14)}; \VUgras{saltano
la corda} \textsuperscript{(13)}, oppure invitano i fratelli e i cugini a una partita
di \VUgras{croquet} \textsuperscript{(15)}. Si divertono un mondo a mandare via con il
loro \VUgras{maglio di legno} la \VUgras{palla} dell'avversario, proprio quando lui
crede di passare comodamente sotto l'archetto!

%%K t02-24-1 p
24. --- I ragazzi preferiscono i giochi più forti. Giocano volentieri ai
\VUgras{birilli} \textsuperscript{(16)}, che abbattono con precisione con la
\VUgras{boccia} \textsuperscript{(17)}. Non disdegnano nemmeno la \VUgras{trottola}
\textsuperscript{(18)} e il \VUgras{bilboquet} \textsuperscript{(19)}, e neanche il
\VUgras{gioco della rana} \textsuperscript{(20)}. Ma i loro giochi preferiti sono le
\VUgras{biglie} \textsuperscript{(21)} e la \VUgras{pallamuro}
\textsuperscript{(22)}, perché il muro della \VUgras{serra} \textsuperscript{(26)} è
perfetto per far rimbalzare la palla leggera.

%%K t02-25-1 p
25. --- Il piccolo Giovanni, in cima ai suoi \VUgras{trampoli} \textsuperscript{(24)},
è molto abile e tiene benissimo l'equilibrio. Vicino a lui alcuni amici giocano a
\VUgras{nascondino} \textsuperscript{(25)} in quella serra e dietro la \VUgras{siepe}.
Altri preferiscono la \VUgras{mano calda} \textsuperscript{(28)} o il
\VUgras{volano} \textsuperscript{(29)}, che vola da una \VUgras{racchetta} all'altra.

%%K t02-noto-1 noto
\VUnotes{143.2mm}{%
(*) I giochi dei bambini hanno spesso nomi puramente nazionali. Li abbiamo nominati
in ido come meglio si poteva, o rifacendoci all'azione stessa che li caratterizza, o
adottando il nome nazionale di un paese o dell'altro quando non è troppo
fuorviante. Per esempio: il \VUgras{gioco della botte} (tedesco e francese) è il
\VUgras{gioco della rana} \textsuperscript{(20)} (inglese \textit{toad in the hole}),
in cui i giocatori cercano di gettare i loro dischi rotondi nella bocca del rospo di
metallo che sta a fauci aperte sulla cassa bucata (la botte). Il
\VUgras{salto sulle spalle} \textsuperscript{(30)} si distingue dalla
\VUgras{cavallina} solo in questo: per saltare al di sopra della testa i giocatori
appoggiano le mani sulle spalle del compagno, mentre nella «\VUgras{cavallina}»
appoggiano le mani solo sulla schiena. --- Oppure ancora: una parte dei giocatori si
china mentre gli altri saltano sopra la loro schiena con le mani sulle spalle; i
primi devono reggere l'urto senza cedere, i secondi devono saltare senza perdere
l'equilibrio (vedi la tavola stessa).%
}

%%K t02-26-1 p
26. --- In mezzo al prato ci sono due alberi. Ai piedi di uno, sette o otto ragazzi
hanno cominciato una partita a \VUgras{salto sulle spalle} \textsuperscript{(30)} (*).
Questo gioco non si conosce in tutti i paesi; ma che cosa sia la \VUgras{cavallina} lo
sanno tutti, e proprio adesso i bambini non ci giocano.

%%K t02-tit-6 sub
\VUsaut{5.0mm}
\VUpk{10.2pt}{40}{Giochi all'aperto.}
\VUsaut{3.4mm}

%%K t02-27-1 p
27. --- Anche la \VUgras{mosca cieca} \textsuperscript{(31)} è molto divertente;
ragazze e ragazzi si mettono a turno la \VUgras{benda} sugli occhi e acchiappano i
più goffi, che si lasciano prendere.

%%K t02-28-1 p
28. --- In mezzo al prato ecco un gioco molto francese, un po' rude per la verità:
l'\VUgras{orso} \textsuperscript{(32)}, molto più rumoroso del tranquillo gioco
dell'\VUgras{aquilone} \textsuperscript{(33)} o perfino del gioco inglese del
\VUgras{tennis} \textsuperscript{(34)}.

%%K t02-29-1 p
29. --- Quest'ultimo sport è la gioia degli alunni più grandi. Perfino i nostri
maestri ci giocano volentieri. Richiede molta abilità e agilità, e a me piace
moltissimo. L'\VUgras{altalena} \textsuperscript{(35)} la lascio alle ragazze.

%%K t02-30-1 p
30. --- Non mi attira nemmeno un altro gioco inglese, che potrebbe facilmente
diventare pericoloso.

%%K t02-30-1b p
Parlo del \VUgras{calcio} \textsuperscript{(36)}, la passione di mio fratello
maggiore. Io preferisco il nostro vecchio \VUgras{gioco delle barre}
\textsuperscript{(38)}, altrettanto sano e molto meno brutale.

%%K t02-tit-7 sub
\VUsaut{4.6mm}
\VUpk{10.2pt}{40}{Ciclismo e giochi con la palla.}
\VUsaut{3.0mm}

%%K t02-31-1 p
31. --- Mi piace anche molto girare intorno al prato in \VUgras{bicicletta}
\textsuperscript{(37)}.

%%K t02-32-1 p
32. --- L'anno prossimo imparerò ad \VUgras{andare a cavallo} \textsuperscript{(39)}.
Un cavallo è uno spettacolo bellissimo quando cammina o trotta con eleganza e supera
gli ostacoli al galoppo!

%%K t02-33-1 p
33. --- D'estate impariamo a \VUgras{nuotare} \textsuperscript{(40)}. Ci tuffiamo e
facciamo il bagno con gioia nell'acqua limpida e fresca del fiume. Qualche volta, alla
fine, facciamo salire un \VUgras{pallone} di carta \textsuperscript{(41)}, che sale
maestoso nell'aria appena si riempie di aria calda.

%%K t02-34-1 p
34. --- Ci sono anche molti altri giochi, ma non finirei più di elencarli, e si capisce
già da questo breve racconto che i giovedì sul nostro prato sono tutt'altro che
noiosi.

%%K t02-34-2 p
N.B. --- \textit{Il maestro potrà facilmente ampliare questo capitolo descrivendo più
in dettaglio ciascuno dei giochi più importanti.}
\VUsaut{6.0mm}
\VUfilet{22mm}
